\subsection{Orthogonal Complements and Minimization Problems}

\begin{exercise}{6c1}
    Suppose $v_1, \ldots, v_m \in V$.
    Prove that
    \[
        \set{v_1, \ldots, v_m}^\perp = (\spn(v_1, \ldots, v_m))^\perp.
    \]
\end{exercise}

\begin{exercise}{6c2}
    Suppose $U$ is a subspace of $V$ with basis $u_1, \ldots, u_m$ and
    \[
        u_1, \ldots, u_m, v_1, \ldots, v_n
    \]
    is a basis of $V$.
    Prove that if the Gram--Schmidt procedure is applied to the basis of $V$ above, producing a list $e_1, \ldots, e_m, f_1, \ldots, f_n$, then $e_1, \ldots, e_m$ is an orthonormal basis of $U$ and $f_1, \ldots, f_n$ is an orthonormal basis of $U^\perp$.
\end{exercise}

\begin{exercise}{6c3}
    Suppose $U$ is the subspace of $\rbb^4$ defined by
    \[
        U = \spn((1, 2, 3, -4),\, (-5, 4, 3, 2)).
    \]
    Find an orthonormal basis of $U$ and an orthonormal basis of $U^\perp$.
\end{exercise}

\begin{exercise}{6c4}
    Suppose $e_1, \ldots, e_n$ is a list of vectors in $V$ with $\norm{e_k} = 1$ for each $k = 1, \ldots, n$ and
    \[
        \norm{v}^2 = |\gen{v, e_1}|^2 + \cdots + |\gen{v, e_n}|^2
    \]
    for all $v \in V$.
    Prove that $e_1, \ldots, e_n$ is an orthonormal basis of $V$.
    \textit{This exercise provides a converse to 6.30(b).}
\end{exercise}

\begin{exercise}{6c5}
    Suppose that $V$ is finite-dimensional and $U$ is a subspace of $V$.
    Show that $P_{U^\perp} = I - P_U$, where $I$ is the identity operator on $V$.
\end{exercise}

\begin{exercise}{6c6}
    Suppose $V$ is finite-dimensional and $T \in \lin{V, W}$.
    Show that
    \[
        T = TP_{(\nul T)^\perp} = P_{\range T}\, T.
    \]
\end{exercise}

\begin{exercise}{6c7}
    Suppose that $X$ and $Y$ are finite-dimensional subspaces of $V$.
    Prove that $P_X P_Y = 0$ if and only if $\gen{x, y} = 0$ for all $x \in X$ and all $y \in Y$.
\end{exercise}

\begin{exercise}{6c8}
    Suppose $U$ is a finite-dimensional subspace of $V$ and $v \in V$.
    Define a linear functional $\phi\colon U \to \fbb$ by
    \[
        \phi(u) = \gen{u, v}
    \]
    for all $u \in U$.
    By the Riesz representation theorem (6.42) as applied to the inner product space $U$, there exists a unique vector $w \in U$ such that
    \[
        \phi(u) = \gen{u, w}
    \]
    for all $u \in U$.
    Show that $w = P_U v$.
\end{exercise}

\begin{exercise}{6c9}
    Suppose $V$ is finite-dimensional.
    Suppose $P \in \lin{V}$ is such that $P^2 = P$ and every vector in $\nul P$ is orthogonal to every vector in $\range P$.
    Prove that there exists a subspace $U$ of $V$ such that $P = P_U$.
\end{exercise}

\begin{exercise}{6c10}
    Suppose $V$ is finite-dimensional and $P \in \lin{V}$ is such that $P^2 = P$ and
    \[
        \norm{Pv} \leq \norm{v}
    \]
    for every $v \in V$.
    Prove that there exists a subspace $U$ of $V$ such that $P = P_U$.
\end{exercise}

\begin{exercise}{6c11}
    Suppose $T \in \lin{V}$ and $U$ is a finite-dimensional subspace of $V$.
    Prove that
    \[
        U \text{ is invariant under } T \iff P_U T P_U = T P_U.
    \]
\end{exercise}

\begin{exercise}{6c12}
    Suppose $V$ is finite-dimensional, $T \in \lin{V}$, and $U$ is a subspace of $V$.
    Prove that
    \[
        U \text{ and } U^\perp \text{ are both invariant under } T \iff P_U T = T P_U.
    \]
\end{exercise}

\begin{exercise}{6c13}
    Suppose $\fbb = \rbb$ and $V$ is finite-dimensional.
    For each $v \in V$, let $\phi_v$ denote the linear functional on $V$ defined by
    \[
        \phi_v(u) = \gen{u, v}
    \]
    for all $u \in V$.
    \begin{subexercises}
        \item Show that $v \mapsto \phi_v$ is an injective linear map from $V$ to $V'$.
        \item Use (a) and a dimension-counting argument to show that $v \mapsto \phi_v$ is an isomorphism from $V$ onto $V'$.
    \end{subexercises}
    \textit{The purpose of this exercise is to give an alternative proof of the Riesz representation theorem (6.42 and 6.58) when $\fbb = \rbb$.
    Thus you should not use the Riesz representation theorem as a tool in your solution.}
\end{exercise}

\begin{exercise}{6c14}
    Suppose that $e_1, \ldots, e_n$ is an orthonormal basis of $V$.
    Explain why the dual basis (see 3.112) of $e_1, \ldots, e_n$ is $e_1, \ldots, e_n$ under the identification of $V'$ with $V$ provided by the Riesz representation theorem (6.58).
\end{exercise}

\begin{exercise}{6c15}
    In $\rbb^4$, let
    \[
        U = \spn((1, 1, 0, 0),\, (1, 1, 1, 2)).
    \]
    Find $u \in U$ such that $\norm{u - (1, 2, 3, 4)}$ is as small as possible.
\end{exercise}

\begin{exercise}{6c16}
    Suppose $C[-1, 1]$ is the vector space of continuous real-valued functions on the interval $[-1, 1]$ with inner product given by
    \[
        \gen{f, g} = \int_{-1}^{1} fg
    \]
    for all $f, g \in C[-1, 1]$.
    Let $U$ be the subspace of $C[-1, 1]$ defined by
    \[
        U = \set{f \in C[-1, 1] \mid f(0) = 0}.
    \]
    \begin{subexercises}
        \item Show that $U^\perp = \set{0}$.
        \item Show that 6.49 and 6.52 do not hold without the finite-dimensional hypothesis.
    \end{subexercises}
\end{exercise}

\begin{exercise}{6c17}
    Find $p \in \pcal_3(\rbb)$ such that $p(0) = 0$, $p'(0) = 0$, and $\int_0^1 |2 + 3x - p(x)|^2\, dx$ is as small as possible.
\end{exercise}

\begin{exercise}{6c18}
    Find $p \in \pcal_5(\rbb)$ that makes $\int_{-\pi}^{\pi} |\sin x - p(x)|^2\, dx$ as small as possible.
    \textit{The polynomial 6.65 is an excellent approximation to the answer to this exercise, but here you are asked to find the exact solution, which involves powers of $\pi$.
    A computer that can perform symbolic integration should help.}
\end{exercise}

\begin{exercise}{6c19}
    Suppose $V$ is finite-dimensional and $P \in \lin{V}$ is an orthogonal projection of $V$ onto some subspace of $V$.
    Prove that $P^\dagger = P$.
\end{exercise}

\begin{exercise}{6c20}
    Suppose $V$ is finite-dimensional and $T \in \lin{V, W}$.
    Show that
    \[
        \nul T^\dagger = (\range T)^\perp \quad \text{and} \quad \range T^\dagger = (\nul T)^\perp.
    \]
\end{exercise}

\begin{exercise}{6c21}
    Suppose $T \in \lin{\fbb^3, \fbb^2}$ is defined by
    \[
        T(a, b, c) = (a + b + c,\, 2b + 3c).
    \]
    \begin{subexercises}
        \item For $(x, y) \in \fbb^2$, find a formula for $T^\dagger(x, y)$.
        \item Verify that the equation $TT^\dagger = P_{\range T}$ from 6.69(b) holds with the formula for $T^\dagger$ obtained in (a).
        \item Verify that the equation $T^\dagger T = P_{(\nul T)^\perp}$ from 6.69(c) holds with the formula for $T^\dagger$ obtained in (a).
    \end{subexercises}
\end{exercise}

\begin{exercise}{6c22}
    Suppose $V$ is finite-dimensional and $T \in \lin{V, W}$.
    Prove that
    \[
        TT^\dagger T = T \quad \text{and} \quad T^\dagger TT^\dagger = T^\dagger.
    \]
    \textit{Both formulas above clearly hold if $T$ is invertible because in that case we can replace $T^\dagger$ with $T^{-1}$.}
\end{exercise}

\begin{exercise}{6c23}
    Suppose $V$ and $W$ are finite-dimensional and $T \in \lin{V, W}$.
    Prove that
    \[
        (T^\dagger)^\dagger = T.
    \]
    \textit{The equation above is analogous to the equation $(T^{-1})^{-1} = T$ that holds if $T$ is invertible.}
\end{exercise}